% !TeX program = pdflatex
\documentclass[12pt]{article}
\usepackage[utf8]{inputenc}
\usepackage{CJKutf8}
\usepackage{amsmath,amssymb,mathtools,amsthm}
\usepackage[a4paper,top=25mm,bottom=25mm,left=20mm,right=20mm]{geometry}
\usepackage{array,booktabs,longtable,tocloft,xcolor}
\usepackage[unicode,colorlinks=true,linkcolor=blue,citecolor=blue,urlcolor=blue]{hyperref}
\usepackage{xurl}
\hypersetup{pdftitle={Direct G-Algebra Kernels and Law-Space Readback}}
\input{glyphtounicode}
\pdfgentounicode=1
\makeatletter
\ifdefined\pdfnobuiltintounicode
  \let\galgebraCJKaddcmap\CJK@addcmap
  \def\CJK@addcmap#1{%
    \ifx\f@encoding\CJK@enc
      \galgebraCJKaddcmap{#1}%
      \pdfnobuiltintounicode\font@name
    \fi}
\fi
\makeatother
\AddToHook{file/c70gkai.fd/after}{%
  \DeclareFontShape{C70}{gkai}{m}{it}{<->CJK*gkaiu}{\CJKnormal}%
  \DeclareFontShape{C70}{gkai}{bx}{it}{<->CJKb*gkaiu}{\CJKbold}%
}
\Urlmuskip=0mu plus 4mu
\setlength{\emergencystretch}{5em}
\setlength{\parindent}{0pt}
\setlength{\parskip}{0.55em}
\renewcommand{\cftsecleader}{\cftdotfill{\cftdotsep}}
\newtheorem{theorem}{定理}[section]
\newtheorem{lemma}[theorem]{引理}
\newtheorem{proposition}[theorem]{命题}
\newtheorem{corollary}[theorem]{推论}
\theoremstyle{definition}
\newtheorem{definition}[theorem]{定义}
\newtheorem{certificate}[theorem]{证书}
\newtheorem{example}[theorem]{例}
\theoremstyle{remark}
\newtheorem{remark}[theorem]{备注}
\newcommand{\QQ}{\mathbb Q}
\newcommand{\RR}{\mathbb R}
\newcommand{\ZZ}{\mathbb Z}
\newcommand{\NN}{\mathbb N}
\newcommand{\CC}{\mathbb C}
\newcommand{\id}{\mathrm{id}}
\newcommand{\lean}[1]{\nolinkurl{#1}}
\newcolumntype{P}[1]{>{\raggedright\arraybackslash}p{#1}}
\DeclareMathOperator{\im}{im}
\DeclareMathOperator{\Reach}{Reach}
\DeclareMathOperator{\Cert}{Cert}
\begin{document}
\begin{CJK*}{UTF8}{gkai}
\renewcommand{\contentsname}{目录}
\renewcommand{\refname}{参考文献}
\renewcommand{\proofname}{证明}
\renewcommand{\abstractname}{摘要}
\renewcommand{\tablename}{表}
\title{G 代数的直接上卷、分型内核与法空间回读\\
\large 经典黎曼猜想与带外力 Navier--Stokes 爆破的回读\\
（G 法空间与 Lean 形式系统）}
\author{高政 （GaoZheng）}
\date{2026-11-25}
\maketitle
\begin{abstract}
本文沿任务商保结构上卷至 G 母体，分型下卷求取见证，以纤维丛协变和多复合反身闭合内核。
全稿登记十四项已证法空间结果，包含七项双向等价、六项充分回读及一项精确读数。覆盖原量词的母体证明经已证回读即为原命题的证明。

经典黎曼分型证明实际 \(\xi\) 的全称矩条件与全部非平凡零点位于临界线等价。
带外力流体分型证明整周期尾方程的多复合闭合，并从合法非零精确周期的存在性
回读零初值、有界动能与光滑紧支外力下的有限时间爆破，保留原黏性量词。
加性数论和全秩自守谱采用各自的等价读数，动理学和宏观流体按精确通量、
共同尺度下的误差条件及有限端点条件验收相应回读。
附录给出有限核验、G 文献定位与 Lean 证明映射。
\end{abstract}

\section*{前言}
\addcontentsline{toc}{section}{前言}
第~\ref{sec:proof-guide} 节登记共同证明链、分型回读及其覆盖范围。
策略格、命中、LHTS、TPH 与 L-路径用于给定有限数据的任务商核验与成本比较
\cite{GNotes,GUnified}。

主丛、纤维作用、高阶相容性采用公开的纯数学构造，法空间的作用与证据采用应用三卷的相应接口
\cite{GPure,GAppI,GAppII,GAppIII}。
任务保持与率—商回读按公开工作笔记和综合文献的相应构造 \cite{GNotes,GUnified}。
附录~\ref{sec:g-internal-basis} 给出本文所用的有限路径、测试、支持因子化和解析输入。
回读的参数、条件及方向在正文定理中登记。

Lean 工程固定为官方 Lean~4.36.0，直接基础模块为 \lean{Init}，版本入口见 \href{https://lean-lang.org/}{Lean 官方站点}。
默认构建核验正文与有限算术附录所用的声明。
附录~\ref{sec:lean-map} 给出正文连续分析、参数化接口与机器核验的对应范围。

\setcounter{tocdepth}{2}
\clearpage
\tableofcontents
\clearpage
\input{g_algebra_rh_ns/parts/proof_guide.texinc}
\input{g_algebra_rh_ns/parts/law.texinc}
\input{g_algebra_rh_ns/parts/chains.texinc}
\input{g_algebra_rh_ns/parts/coverage.texinc}
\input{g_algebra_rh_ns/parts/repair.texinc}
\input{g_algebra_rh_ns/parts/targets.texinc}
\appendix
\input{g_algebra_rh_ns/parts/internal_basis.texinc}
\section{有限算术核验}
\label{sec:finite-zeta-verification}
有理包围、轮廓计数及变号证书的结论按各自给定的有限域读取。
\input{g_algebra_rh_ns/parts/arithmetic_enclosures.texinc}
\input{g_algebra_rh_ns/parts/zeta_window.texinc}
\input{g_algebra_rh_ns/parts/mapping.texinc}
\end{CJK*}
\end{document}
